\documentclass[12pt]{article}
\usepackage[utf8]{inputenc}
\usepackage[T1]{fontenc}
\usepackage{hyperref}
\title{Malgana Language Resources for K-12 Curriculum Development}
\author{Compiled by Claude Sonnet 5 (claude-sonnet-5) for Andr\'as Kornai}
\date{}
\begin{document}
\maketitle

\section{Introduction}
Malgana (ISO 639-3 vml; Glottocode malg1242) is a Pama-Nyungan language of the Shark Bay region of Western Australia. It has no fluent native speakers, and revival work draws on archival recordings and written sources.

\section{Primary Educational Resources}
Search results indicate that vocabulary is taught in the local school, and that school and community language classes are mentioned on a heritage page. No published textbook or primer was found.
\begin{itemize}
\item Shark Bay community page on Aboriginal languages: \url{https://www.sharkbay.org/?p=165}
\item Gutharraguda Cultural Heritage page: \url{https://www.sharkbayvisit.com.au/gutharraguda-cultural-heritage/}
\end{itemize}

\section{Linguistic and Reference Materials}
\begin{itemize}
\item Gargett, \emph{A Salvage Grammar of Malgana} (2011), ANU repository PDF: \url{https://openresearch-repository.anu.edu.au/bitstream/1885/146771/1/Gargett-2011-SalvageGrammarMalgana.pdf}
\item Same work, Open University repository record: \url{https://oro.open.ac.uk/73150}
\item Wikipedia article on the Malgana language. It mentions an illustrated wordlist produced by the Irra Wangga Language Centre: \url{https://en.wikipedia.org/wiki/Malgana_language}
\end{itemize}

\section{Religious and Community Resources}
Community heritage pages describe Elders' storytelling sessions and partnerships with language centres (see the Gutharraguda page above). The sources disagree on which language centre is currently involved, so check this directly.

\section{Digital and Audio Resources}
Grammar work reportedly draws on 1990s audio recordings. No public audio collection was found in search results.

\section{Curriculum Development Context}
A heritage and revitalization curriculum would build on the salvage grammar, the illustrated wordlist, and local school and community programmes. Contact details were not verified here. Whether the wordlist is publicly downloadable was not checked.
\end{document}